% Source: https://artofproblemsolving.com/community/c3h2598208_interesting_paradox
% Stable identity: aops-topic-2598208-q01
\textbf{(a)} Stage $k$ covers $\dfrac1{2^k}$ of the distance, so it takes $\dfrac{60}{2^k}$ minutes. The first $n$ stages take
\[60\left(\dfrac12+\dfrac14+\cdots+\dfrac1{2^n}\right)=60\left(1-\dfrac1{2^n}\right)\]
minutes, by the geometric series formula (or because after $n$ stages a fraction $\dfrac1{2^n}$ of the distance is left).

\textbf{(b)} Infinitely many positive numbers can have a finite sum when they shrink fast enough. Here the total time after $n$ stages is always less than $60$ and gets as close to $60$ as we like, so the infinite sum is
\[\dfrac{60}2+\dfrac{60}4+\dfrac{60}8+\cdots=\boxed{60\text{ minutes}},\]
the time of the whole walk, as it must be.
